% A \noexpand'ED TOKEN IS ITSELF IF IT CANNOT BE EXPANDED.
%
% \noexpand does not hide what it marks. What comes back is:
%
%   the token itself          when its meaning cannot be expanded
%                             -- a \chardef, a register, a primitive, \relax
%   something \relax-like     when it can -- a macro, an undefined name --
%   but NOT \relax            which is \ifx-different from \relax as well
%
% so `\expandafter\ifx\noexpand\x\x' is TRUE exactly when \x is
% unexpandable. HSTeX kept a mark of its own on every \noexpand'ed control
% sequence and \ifx saw the mark rather than the meaning, so the test was
% false for everything.
%
% expl3 turns on this test: it is how a V-expansion tells one of its
% variables from a macro. Without it \seq_gpush:NV pushed the name of a
% \chardef'd read stream instead of its number, so every \ior_open after
% the first drew "! Bad number (16)" -- 16 being the terminal stream that
% a closed ior is \let back to -- and every LaTeX format load reported it.
%
% \ifx reads two tokens and skips no spaces, which is why the last case
% compares a letter against the space in front of it and is false.
\catcode`\{=1 \catcode`\}=2
\tracingonline=1
\chardef\ca=3
\countdef\co=5
\def\ma{x}
\let\re=\relax
\message{[A a chardef token]}
\expandafter\ifx\noexpand\ca\ca\message{[same]}\else\message{[diff]}\fi
\expandafter\ifx\noexpand\ca\relax\message{[relax]}\else\message{[notrelax]}\fi
\message{[B a countdef token]}
\expandafter\ifx\noexpand\co\co\message{[same]}\else\message{[diff]}\fi
\message{[C a macro]}
\expandafter\ifx\noexpand\ma\ma\message{[same]}\else\message{[diff]}\fi
\expandafter\ifx\noexpand\ma\relax\message{[relax]}\else\message{[notrelax]}\fi
\message{[D relax itself]}
\expandafter\ifx\noexpand\re\re\message{[same]}\else\message{[diff]}\fi
\message{[E a primitive]}
\expandafter\ifx\noexpand\hbox\hbox\message{[same]}\else\message{[diff]}\fi
\message{[F an undefined cs]}
\expandafter\ifx\noexpand\zz\zz\message{[same]}\else\message{[diff]}\fi
\expandafter\ifx\noexpand\zz\relax\message{[relax]}\else\message{[notrelax]}\fi
\message{[G a cs let to a primitive]}
\let\bx=\hbox
\expandafter\ifx\noexpand\bx\hbox\message{[same]}\else\message{[diff]}\fi
\message{[H what noexpand does to a character token]}
\expandafter\ifx\noexpand A A\message{[same]}\else\message{[diff]}\fi
\message{[done]}
\end
